\chapter{Resultados e Discussão}
\label{capitulo:resultados_discussao}

Este capítulo apresenta e debate os resultados obtidos após a implementação do simulador \textit{Royal Legacy}. São expostas as interfaces gráficas finais, traçando um comparativo com os protótipos idealizados, além da apresentação das métricas de desempenho que validam a eficiência da arquitetura assíncrona e do sistema de temas dinâmicos.

\section{Apresentação das interfaces finais}
\label{secao:interfaces_finais}

A interface gráfica foi refinada na Godot Engine com a aplicação de uma estética \textit{pixel art} coerente e paletas de cores harmônicas. O resultado final priorizou a clareza visual, essencial para um jogo de tabuleiro.

\subsection{Menus e Customização}

A Figura \ref{figura:resultado_menu} apresenta a interface de seleção de temas e configurações de áudio. Diferente do rascunho inicial, a implementação final utiliza o \textit{script} \texttt{tela\_temas.gd} para alternar instantaneamente entre paletas (como o tema "Clássico" e "Madeira"). Além disso, o menu de sons (\texttt{tela\_sons.gd}) permite o controle granular dos decibéis para os canais \textit{Master}, Música e Efeitos, garantindo uma experiência auditiva personalizada.

\figura{./capitulos/resultados_discussao/tela_menu_final.png}{12cm}{Interface final de seleção de temas e controles de áudio.}{figura:resultado_menu}

\subsection{Tabuleiro e Gameplay}

A Figura \ref{figura:resultado_jogo} exibe a tela de jogo principal. O refinamento visual incluiu a renderização de peças personalizadas e a implementação de sinais visuais de seleção. O uso da classe \texttt{base\_xadrez.gd} permitiu que a lógica de movimentação respondesse de forma fluida, com o realce das casas de origem e destino, facilitando a compreensão do fluxo da partida pelo usuário.

\figura{./capitulos/resultados_discussao/tela_jogo_final.png}{12cm}{Interface de gameplay com aplicação de temas e realce de movimentos.}{figura:resultado_jogo}

\section{Avaliação de desempenho e estabilidade}
\label{secao:avaliacao_desempenho}

Para validar a solução de integração assíncrona, o sistema foi testado monitorando-se a taxa de quadros (\textit{Frames Per Second} - FPS) da Godot enquanto o \textit{middleware} Python processava as jogadas da IA. Os testes foram divididos nos três níveis de dificuldade programados em \texttt{chess\_bridge.py}.

Os resultados, consolidados na Tabela \ref{tabela:resultados_desempenho}, demonstram o tempo médio entre o envio da string FEN e o recebimento da jogada UCI, sob uma taxa de atualização constante de 60 FPS.

\begin{table}[htb]
	\centering
	\caption{Métricas de tempo de resposta da IA e estabilidade da interface.}
	\label{tabela:resultados_desempenho}
	\begin{tabular}{|c|c|c|c|}
		\hline
		\textbf{Dificuldade} & \textbf{Lógica da IA} & \textbf{Tempo de Resposta} & \textbf{Estabilidade (FPS)} \\ \hline
		Fácil & Movimento Aleatório & 0,50 segundos* & Constante a 60 FPS \\ \hline
		Médio & Stockfish (Depth 1) & 0,85 segundos & Constante a 60 FPS \\ \hline
		Difícil & Stockfish (Depth 20) & 2,10 segundos & Constante a 60 FPS \\ \hline
	\end{tabular}
\vspace{2pt}
{\small \\ Fonte: Elaborado pelo autor (2026). *Inclui o atraso proposital de \textit{sleep} para simulação de pensamento.}
\end{table}

\section{Discussão dos resultados}
\label{secao:discussao}

A análise dos dados confirma que o principal problema técnico — o congelamento da interface durante o processamento da IA — foi mitigado com sucesso. O uso de \texttt{ia\_thread.start()} no \textit{script} \texttt{base\_xadrez.gd} permitiu que, mesmo quando o Stockfish operava em profundidade máxima (nível Difícil), o motor gráfico da Godot continuasse processando entradas do usuário e animações de interface sem interrupções.

No nível "Fácil", o tempo de resposta de 0,5 segundos foi mantido por uma pausa programada (\textit{time.sleep}) no Python, o que humaniza a experiência ao evitar respostas instantâneas que poderiam confundir o jogador iniciante. Já nos níveis superiores, o \textit{middleware} provou ser resiliente ao gerenciar a comunicação via \textit{pipes} de terminal sem causar vazamento de memória ou picos de processamento na via principal (\textit{main thread}).

Por fim, a modularização do projeto em cenas independentes (\texttt{estruturador.gd}) mostrou-se eficiente, mantendo o consumo de memória RAM estável, uma vez que elementos pesados como o tabuleiro só são instanciados após a seleção do modo de jogo, validando as escolhas de arquitetura de \textit{software} definidas no Capítulo \ref{capitulo:projeto}.